% Generated draft: requires mathematical and source review.
\documentclass[11pt]{article}
\usepackage[margin=25mm]{geometry}
\usepackage{fontspec}
\setmainfont{Latin Modern Roman}
\usepackage{amsmath,amssymb,mathtools,unicode-math}
\setmathfont{Latin Modern Math}
\usepackage{enumitem,needspace,xurl,hyperref}
\setlength{\parindent}{0pt}
\setlength{\parskip}{5pt}
\setlength{\emergencystretch}{3em}
\newcommand{\N}{\mathbb{N}}
\newcommand{\Z}{\mathbb{Z}}
\newcommand{\Q}{\mathbb{Q}}
\newcommand{\R}{\mathbb{R}}
\newcommand{\C}{\mathbb{C}}
\newcommand{\F}{\mathbb{F}}
\newcommand{\A}{\mathbb{A}}
\newcommand{\PP}{\mathbb{P}}
\newcommand{\E}{\mathbb{E}}
\newcommand{\eps}{\varepsilon}
\newcommand{\dd}{\,\mathrm{d}}
\begin{document}
\section*{UnsolvedMath ID: 800924}
\textbf{Source problem code:} catalogue entry
\subsection{Definitions and mathematical statement}
For a smooth Riemannian manifold \((M,g)\) without boundary, write \(\operatorname{Vol}_g(M)=\int_M \dd\operatorname{vol}_g\) for its Riemannian volume. Completeness means geodesic completeness: every maximal geodesic is defined on all of \(\R\). A closed geodesic means a nonconstant smooth periodic geodesic, equivalently a nonconstant smooth map \(\gamma:\R/T\Z\to M\), for some \(T>0\), satisfying \(\nabla_{\dot\gamma}\dot\gamma=0\), where \(\nabla\) is the Levi--Civita connection of \(g\).

\textbf{Question (attributed to Bangert).} Does every complete Riemannian manifold \((M,g)\) of finite volume admit a closed geodesic? In the boundaryless interpretation above, the requested implication is
\[
\bigl((M,g)\text{ is geodesically complete and }\operatorname{Vol}_g(M)<\infty\bigr)
\quad\Longrightarrow\quad
\exists T>0\ \exists\gamma:\R/T\Z\to M\quad
\bigl(\gamma\text{ is smooth and nonconstant, }\nabla_{\dot\gamma}\dot\gamma=0\bigr).
\]
The supplied statement specifies neither a dimension range nor connectedness, nonemptiness, or boundary conventions. The definitions above explicitly use the smooth boundaryless interpretation; the intended scope, particularly any positive-dimensional restriction, requires confirmation. No compactness assumption is supplied.
\subsection{Short English statement}
The question attributed to Bangert asks whether completeness and finite Riemannian volume force the existence of a nonconstant periodic geodesic. The intended dimension range and manifold conventions require clarification. The supplied registry labels the problem open; this is not an independently verified status assessment.
\subsection{Sources}
\begin{itemize}
\item Local UnsolvedMath registry, permanent ID \(800924\), exact source problem code \texttt{catalogue entry}, title \emph{Open Problems and Questions about Geodesics --- Question 4}. The supplied record lists Burns and Matveev as proposers and gives the year 2019; the question itself is attributed to Bangert.
\item Burns and Matveev, \emph{Open Problems and Questions about Geodesics} (2019, as dated in the supplied record), Question 4. Supplied source URL: \url{https://arxiv.org/abs/1308.5417}. The supplied provenance records access on 2026-10-16 and extraction by \texttt{arxiv\_latex}; these are supplied provenance claims, not independently verified access or verification by this draft.
\end{itemize}
\end{document}
